\documentclass[11pt]{article}
\renewcommand{\labelitemi}{--}

\begin{document}

\title{The FFmpeg Codebase: A Guided Tour}
\author{C libraries, pipes, packets and frames}
\date{November 2026}
\maketitle

\section{The one-sentence version}

FFmpeg is not really one program. It is a family of C libraries that implement
a media pipeline, plus a few small command-line programs --- \verb|ffmpeg|,
\verb|ffprobe| and \verb|ffplay| --- that wire those libraries together.

\begin{verbatim}
file / URL / camera
      |
      v  libavformat     demux: identify the container, split out streams
   AVPacket              compressed bytes
      |
      v  libavcodec      decode
   AVFrame               raw pixels / raw samples
      |
      v  libavfilter     scale, crop, overlay, volume, subtitles
   AVFrame
      |
      v  libavcodec      encode
   AVPacket
      |
      v  libavformat     mux: write MP4 / MKV / TS
   output file
\end{verbatim}

The usual early confusion: demuxing is not decoding, and encoding is not
muxing. A demuxer pulls compressed packets out of a container; a decoder turns
packets into raw frames. The encoder and muxer do the reverse.

\section{Repository layout}

\bigskip

\begin{center}
\begin{tabular}{@{}p{0.26\textwidth} p{0.66\textwidth}@{}}
\hline
\textbf{Directory} & \textbf{What lives there} \\
\hline
\verb|libavutil/|    & Foundation: \verb|AVFrame|, \verb|AVBufferRef|, \verb|AVRational|, formats, math, logging \\
\verb|libavcodec/|   & Biggest directory. Decoders/encoders: \verb|h264dec.c|, \verb|hevc/|, \verb|libx264.c|, \verb|libx265.c|, bitstream readers, SIMD in \verb|x86/| \\
\verb|libavformat/|  & Demuxers/muxers and protocols: \verb|mov.c|, \verb|movenc.c|, \verb|matroskadec.c|, \verb|mpegts.c|, \verb|http.c| \\
\verb|libavfilter/|  & Filter graph engine plus filters: \verb|vf_scale.c|, \verb|vf_crop.c|, \verb|af_volume.c| \\
\verb|libswscale/|   & Pixel-format conversion and resizing: \verb|sws_scale()| \\
\verb|libswresample/|& Audio conversion: \verb|swr_convert()|, sample rate, channel layout \\
\verb|libavdevice/|  & Capture/output backends: \verb|x11grab.c|, \verb|v4l2.c|, \verb|pulse.c| \\
\verb|fftools/|      & The CLI programs: \verb|ffmpeg.c|, \verb|ffprobe.c|, \verb|ffplay.c|, plus split-out demux/decode/filter/encode/mux files \\
\verb|configure|     & Custom configure script; generates \verb|config.h| and \verb|config.mak| \\
\verb|ffbuild/|      & Make fragments used by the build \\
\verb|doc/examples/| & Small complete example programs --- the best entry point \\
\verb|tests/|        & FATE regression suite \\
\hline
\end{tabular}
\end{center}

\section{The types you must know}

\bigskip

\begin{center}
\begin{tabular}{@{}p{0.26\textwidth} p{0.66\textwidth}@{}}
\hline
\textbf{Type} & \textbf{Meaning} \\
\hline
\verb|AVFormatContext| & One open input or output. Holds streams, duration, metadata, I/O context \\
\verb|AVStream|        & One track inside the file: codec id, size, sample rate, \verb|time_base| \\
\verb|AVCodecContext|  & Live codec state: negotiated dimensions, pixel format, quality, private options \\
\verb|AVPacket|        & Compressed chunk: \verb|data|, \verb|size|, \verb|pts|, \verb|dts|, \verb|stream_index| \\
\verb|AVFrame|         & Decoded raw data: video planes or audio samples, plus \verb|pts| and format \\
\hline
\end{tabular}
\end{center}

Two supporting ideas matter more than they look:

\begin{itemize}
  \item \verb|AVRational time_base|: timestamps are integers interpreted in a
        per-stream fraction such as \verb|{1, 90000}|. Many A/V sync bugs are
        timebase bugs.
  \item \verb|AVBufferRef|: reference-counted memory. Packets and frames rarely
        copy pixels; they pass refs around. Mutating a frame you do not own is
        a classic mistake.
\end{itemize}

\section{Read the pipeline in code}

The skeleton of \verb|doc/examples/demux_decode.c| looks like this:

\begin{verbatim}
avformat_open_input(&fmt, "in.mkv", NULL, NULL);
avformat_find_stream_info(fmt, NULL);

for (each stream) {
    const AVCodec *dec = avcodec_find_decoder(par->codec_id);
    avcodec_alloc_context3(dec);
    avcodec_parameters_to_context(dec_ctx, par);
    avcodec_open2(dec_ctx, dec, NULL);
}

while (av_read_frame(fmt, &pkt) >= 0) {
    avcodec_send_packet(dec_ctx, &pkt);
    while (avcodec_receive_frame(dec_ctx, frame) == 0) {
        /* AVFrame: raw pixels/samples -- filters, then encoder */
    }
}
\end{verbatim}

Modern FFmpeg uses the send/receive API. Old tutorials that call
\verb|avcodec_decode_video2()| describe an API removed years ago.

\section{How one codec is implemented}

Take H.264 as an example:

\begin{itemize}
  \item \verb|AVCodec| is a struct of function pointers: \verb|init|,
        \verb|decode|, \verb|close|, \verb|flush|, and capability flags.
  \item \verb|AVCodecContext->priv_data| points to a codec-specific struct
        holding SPS/PPS, slice state and reference frames.
  \item Bitstreams are read with \verb|get_bits.h| helpers such as
        \verb|get_bits()| and \verb|show_bits()|, not byte-by-byte string logic.
  \item Shared H.264/H.265 syntax lives in \verb|h2645_parse.c|; H.265 has its
        own \verb|libavcodec/hevc/| directory.
  \item Hot loops have hand-written SIMD in \verb|libavcodec/x86/| and
        \verb|arm/| --- hence the \verb|_sse2|, \verb|_avx2|, \verb|_neon|
        variants chosen at runtime.
  \item Third-party encoders are wrappers: \verb|libx264.c|, \verb|libx265.c|,
        \verb|libvpxenc.c|, \verb|libaomenc.c|. They exist only if configure
        found the external library.
\end{itemize}

That last point explains the VLC codec problem from earlier. Fedora's
\verb|ffmpeg-free| was configured with
\verb|--disable-decoder='h264,hevc,vc1,vvc'|, so \verb|CONFIG_HEVC_DECODER| was
zero and those \verb|AVCodec| entries were never compiled in. VLC's
\verb|libavcodec_plugin.so| had no decoder for \verb|avcodec_find_decoder()| to
find. Installing the RPM Fusion build flipped those flags on.

\section{The build system}

\verb|./configure| probes the compiler and libraries, then generates
\verb|config.h|, \verb|config.mak| and \verb|ffbuild/config.mak|. Source files
are full of guards such as \verb|#if CONFIG_HEVC_DECODER|, and Makefiles pull
in objects conditionally. That is why flags like \verb|--enable-gpl|,
\verb|--enable-libx265| and \verb|--enable-cuda| change the resulting binary.

Tests are FATE: reference clips plus expected hashes, described in
\verb|tests/fate/*.mak| and run with \verb|make fate|.

\section{Best reading order}

\begin{enumerate}
  \item \verb|doc/examples/demux_decode.c| and \verb|doc/examples/transcode.c|:
        small, complete, no CLI noise.
  \item \verb|libavformat/demux.c|: where \verb|av_read_frame()| lives now.
  \item \verb|libavcodec/decode.c| and \verb|encode.c|: send/receive machinery
        and threading.
  \item \verb|libavfilter/graphparser.c|: how a string such as
        \verb|scale=1280:-1,eq=gamma=1.1| becomes a graph of filter nodes.
  \item Only then \verb|fftools/ffmpeg.c|, which is the full conductor and is
        harder to read cold, especially after the recent split into
        \verb|ffmpeg_demux.c|, \verb|ffmpeg_enc.c| and \verb|ffmpeg_mux.c|.
\end{enumerate}

For exploration without reading C: \verb|ffprobe -show_streams file.mkv|,
\verb|ffmpeg -decoders|, \verb|ffmpeg -h encoder=libx265|, \verb|ffmpeg -filters|
and \verb|ffmpeg -loglevel debug| are all useful.

\end{document}
